\chapter{Conclusiones}
\label{ch:conclusiones}
% No es un resumen: interpreta. Una conclusión por objetivo. Se cierra tras Resultados.

\section{Cumplimiento de objetivos}
\label{sec:concl-objetivos}
% Repasar H1-H6; un resultado negativo robusto (H1 sí, H2 no) es contribución válida.
% Aquí vive la interpretación de los datos y el contraste con el estado del arte (\ref{ch:estado-del-arte});
% Resultados solo describe (guía ETSINF). Una conclusión parcial por objetivo.

\section{Aportaciones}
\label{sec:concl-aportaciones}
% Comparativa multi-familia, barrido temporal, coste asintótico frente a capacidad predictiva.

\section{Limitaciones}
\label{sec:concl-limitaciones}
% Solo visión, correlacional, sin augmentation. El barrido por ejemplo recorre todos los parámetros en las tres redes; solo GWA lee la última capa.
% Las líneas de continuación viven ahora en su propio capítulo, \ref{ch:trabajos-futuros}.
